% Chromatographie sur couche mince : cuve et plaque révélée (L'essentiel, chapitre 9)
\begin{tikzpicture}[x=1cm,y=1cm]
% Cuve
\filldraw[fill=white,draw=otNavy!70,line width=0.7pt] (-1.1,0) rectangle (1.1,3.3);
\filldraw[pcAppareil] (-1.2,3.3) rectangle (1.2,3.45);
\fill[pcEau] (-1.05,0.05) rectangle (1.05,0.45);
\filldraw[fill=white,draw=otInk!55,line width=0.5pt] (-0.55,0.15) rectangle (0.55,3.1);
\draw[otMuted,line width=0.4pt,dash pattern=on 1pt off 1pt] (-0.55,0.7) -- (0.55,0.7);
\foreach \x in {-0.3,0,0.3}{\fill[otNavy!60] (\x,0.7) circle (0.05);}
\pcNote[west]{(1.05,0.3)}{(1.4,0.5)}{éluant}
\pcNote[west]{(0.55,0.7)}{(1.4,1.2)}{ligne de dépôt}
\node[pcLegendeG,anchor=north] at (0,-0.15) {élution dans la cuve};
% Plaque révélée
\begin{scope}[shift={(4.6,0)}]
  \filldraw[fill=white,draw=otInk!55,line width=0.5pt] (-0.9,0) rectangle (0.9,3.3);
  \draw[otMuted,line width=0.4pt,dash pattern=on 1pt off 1pt] (-0.9,0.5) -- (0.9,0.5);
  \draw[otGreenDark,line width=0.5pt,dash pattern=on 2pt off 1pt] (-0.9,2.9) -- (0.9,2.9);
  \foreach \x/\l in {-0.55/A,0/B,0.55/R}{\fill[otMuted] (\x,0.5) circle (0.035); \node[pcLegendeG,anchor=north] at (\x,0.42) {\l};}
  \fill[otNavy!70] (-0.55,1.7) ellipse (0.13 and 0.09) (0,1.7) ellipse (0.13 and 0.09) (0.55,1.7) ellipse (0.13 and 0.09);
  \fill[otOrange!80] (-0.55,2.42) ellipse (0.11 and 0.08);
  \pcNote[west]{(0.9,2.9)}{(1.3,3.05)}{front de l'éluant}
  \pcNote[west]{(0.68,1.7)}{(1.3,1.75)}{produit}
  \pcNote[west]{(-0.44,2.42)}{(1.3,2.45)}{impureté}
  \node[pcLegendeG,anchor=north] at (0,-0.15) {plaque révélée};
\end{scope}
\end{tikzpicture}
\par\smallskip\legende{A : produit brut ; B : produit purifié ; R : référence. Le produit purifié ne donne qu'une tache, à la même hauteur que la référence.}
